\section{引言}

从一段由多个镜头组成的单目视频中恢复全局一致的多人体运动，需要在相机视角突然变化时连接不同镜头中的人体轨迹
~\citep{pavlakos2022multishot,zhang2025humanmm,on2026multithumbs}。
镜头切换会打破相机运动的连续性，并显著改变人体外观，使各镜头的重建落入彼此不兼容的世界坐标系，人物身份也可能随之断裂。

本文中的“多镜头”是指单目视频中按时间依次出现的镜头片段。\method{} 每次只接收一帧，从不联合观察同一时刻的同步多视角输入。

已有多镜头方法通常联合处理相邻镜头，或对已经观测到的完整序列进行优化
~\citep{zhang2025humanmm,on2026multithumbs}，因而无法在帧到达时递归地产生结果。流式重建按顺序处理输入，用固定大小的状态维护历史，并避免逐视频全局优化；然而，现有方法通常假设观测在时间上连续
~\citep{wang2025cut3r,chen2026human3r,zhao2026onlinehmr}。

镜头切换由此暴露出递归历史相互冲突的作用：历史信息为镜头间对齐提供参照，但继续传播同一状态可能引起镜头内相机漂移；重新初始化避免了这种状态延续，却会失去跨镜头参照。我们的核心观察是，对齐与后续递归应被视为两个独立决策：旧状态可以临时用于估计镜头间关系，独立重新初始化的状态则负责新镜头中的重建。同一检查点下的受控实验进一步表明，跨镜头延续状态会在新镜头内部引入随时间变化的相机漂移，而不仅是一个固定坐标偏移（第~\ref{sec:ablation-zh} 节）。

基于这一观察，我们提出 \method{}。在镜头边界，历史条件对齐路径使用旧状态和当前图像估计镜头间变换，随后丢弃其临时状态；独立重新初始化的路径产生唯一继续传播的状态。基于预测几何的人物关联连接跨镜头身份，一个共享变换将相机与人体共同映射到已有世界坐标系。整个过程不访问未来帧、不修改已输出结果，也不执行逐视频优化。

多个多人体基准上的实验表明，\method{} 在保持流式推理的同时改善了世界坐标人体重建、相机一致性和身份连续性，并在大视角变化下获得更明显的优势。

本文的贡献如下：

\begin{itemize}
  \item 我们识别并通过实验刻画了递归历史在镜头切换处的冲突作用：切换前历史为镜头间对齐提供参照，继续传播同一状态可能引起镜头内相机漂移，而单独重新初始化会失去跨镜头参照。
  \item 我们通过临时的历史条件对齐路径与独立重新初始化的递归路径实现这种解耦，并以基于几何的人物关联和相机—人体共享变换维持空间与身份一致性。
  \item 在多个多人体基准上，我们改善了世界坐标人体重建、相机一致性和身份连续性；在大视角变化下优势更加明显，且无需未来帧或逐视频优化。
\end{itemize}
